\chapter{Fundamentação Teórica}
\label{cap:fundamentacao-teorica}

Este capítulo apresenta os fundamentos teóricos que embasam a pesquisa, estruturados em
cinco eixos: os Objetivos de Desenvolvimento Sustentável e sua mensuração em escala local; o
Processamento de Linguagem Natural; a análise de sentimentos; a arquitetura Transformer; e os
Modelos de Linguagem de Grande Porte, incluindo o \gls{BERT} e suas variantes.

\section{Objetivos de Desenvolvimento Sustentável}
\label{sec:ods}

Os Objetivos de Desenvolvimento Sustentável (\gls{ODS}) integram a Agenda 2030 da
Organização das Nações Unidas, adotada em setembro de 2015 por 193 países membros
\cite{onu2015}. Compostos por 17 objetivos e 169 metas, os ODS constituem um arcabouço global
de orientação para políticas públicas, investimentos e ações coletivas voltadas ao desenvolvimento
sustentável nas dimensões econômica, social e ambiental. Diferentemente dos Objetivos de
Desenvolvimento do Milênio (ODM), os ODS abrangem tanto países em desenvolvimento quanto
desenvolvidos e reconhecem explicitamente a interdependência entre prosperidade, equidade e
preservação ambiental.

No Brasil, o acompanhamento dos ODS em escala municipal é realizado, entre outros
instrumentos, pelo Índice de Desenvolvimento Sustentável das Cidades (\gls{IDSC}-BR)
\cite{idsc2024}, que consolida indicadores primários de fontes como o IBGE, registros
administrativos e pesquisas oficiais para mais de cinco mil municípios. O IDSC-BR classifica o
desempenho dos municípios em faixas de 0 a 100 por ODS, permitindo identificar prioridades de
política pública. Embora robusto para análises comparativas nacionais, o índice apresenta
limitações de granularidade espacial e defasagem temporal, características que esta dissertação
busca complementar.

\subsection{ODS 9 -- Indústria, Inovação e Infraestrutura}
\label{subsec:ods9}

O ODS 9 tem como propósito construir infraestruturas resilientes, promover a
industrialização inclusiva e sustentável e fomentar a inovação \cite{onu2015}. Suas metas cobrem
desde a ampliação de infraestruturas de transporte e logística acessíveis à população de baixa
renda até o estímulo à pesquisa científica, ao acesso universal a tecnologias de informação e
comunicação e ao suporte a micro e pequenas empresas nos países em desenvolvimento.

Em municípios brasileiros de médio porte, o desempenho no ODS 9 é frequentemente
comprometido pela degradação de infraestruturas de transporte público (terminais rodoviários,
terminais urbanos), pela baixa conectividade digital em periferias e pela ausência de ecossistemas
de inovação que apoiem o empreendedorismo local. Mossoró-RN apresenta pontuação muito baixa
(faixa 0--39,99) nesse objetivo \cite{idsc2024}, sem que os indicadores agregados revelem os
fatores específicos que explicam esse desempenho.

\subsection{ODS 16 -- Paz, Justiça e Instituições Eficazes}
\label{subsec:ods16}

O ODS 16 busca promover sociedades pacíficas e inclusivas, proporcionar o acesso à
justiça para todos e construir instituições eficazes, responsáveis e inclusivas em todos os níveis
\cite{onu2015}. Suas metas abrangem a redução da violência e das mortes relacionadas a conflitos,
o combate à corrupção e ao suborno, o acesso equitativo à justiça, a proteção das liberdades
fundamentais e a garantia de representação nos processos decisórios.

No contexto municipal, o ODS 16 se manifesta na qualidade dos serviços prestados por
órgãos públicos (cartórios, delegacias, varas judiciais, serviços de assistência social) e na
percepção cidadã sobre eficiência institucional e segurança. A análise de avaliações públicas de
estabelecimentos vinculados a esses serviços permite capturar dimensões de desempenho que
raramente aparecem em indicadores oficiais agregados.

\subsection{ODS 3 -- Saúde e Bem-Estar}
\label{subsec:ods3}

O ODS 3 objetiva assegurar uma vida saudável e promover o bem-estar para todos, em
todas as idades \cite{onu2015}. Suas metas incluem a redução da mortalidade materna e infantil,
o combate a doenças transmissíveis e não transmissíveis, o fortalecimento da cobertura universal
de saúde e o acesso a medicamentos e vacinas essenciais. Em escala municipal, traduz-se na
qualidade dos serviços de atenção básica, na disponibilidade de leitos hospitalares, no tempo de
espera em unidades de saúde e na percepção dos usuários sobre o atendimento recebido.

\subsection{ODS 4 -- Educação de Qualidade}
\label{subsec:ods4}

O ODS 4 visa garantir educação inclusiva, equitativa e de qualidade, promovendo
oportunidades de aprendizagem ao longo da vida para todos \cite{onu2015}. Em nível municipal,
abrange a qualidade do ensino nas redes pública e privada, a infraestrutura escolar, a taxa de
conclusão da educação básica e o acesso à educação profissional e superior. Avaliações públicas
de estabelecimentos educacionais permitem capturar percepções sobre ambiente físico, qualidade
do ensino, relação professor-aluno e gestão escolar, complementando os dados do Índice de
Desenvolvimento da Educação Básica (IDEB).

\section{Mensuração Local dos ODS}
\label{sec:mensuracao-local}

A mensuração do progresso nos ODS em escala subnacional apresenta desafios
metodológicos e operacionais relevantes \cite{abidoye2022}. Indicadores nacionais ou estaduais
tendem a obscurecer disparidades locais críticas: um município pode apresentar desempenho
médio satisfatório enquanto bairros ou setores específicos enfrentam privações severas. A
localização dos ODS --- entendida como a adaptação do monitoramento à escala municipal e
intraurbana --- requer fontes de dados granulares e metodologias capazes de capturar variação
espacial.

Nesse contexto, dados digitais produzidos por cidadãos em plataformas \textit{online} emergem
como fontes complementares promissoras \cite{fu2023,larosa2025}. Avaliações públicas
georreferenciadas, como as disponibilizadas pela Google Places API \cite{googleplaces2024},
vinculam relatos textuais espontâneos a estabelecimentos físicos concretos --- terminais
rodoviários, unidades de saúde, escolas, órgãos públicos, espaços urbanos. Essa ancoragem
geoespacial possibilita diagnósticos em nível de estabelecimento, muito mais específicos do que
os indicadores agregados do IDSC-BR, e permite identificar \emph{onde} e \emph{por que} os
problemas se manifestam.

A abordagem perceptiva complementa --- sem substituir --- os indicadores oficiais. Enquanto
o IDSC-BR oferece comparabilidade temporal e cobertura nacional, os dados do Google Places
fornecem especificidade local, atualidade (avaliações são inseridas continuamente) e riqueza
qualitativa que os indicadores quantitativos não conseguem capturar \cite{wankhade2022}.

\section{Processamento de Linguagem Natural}
\label{sec:pln}

O Processamento de Linguagem Natural (\gls{NLP}) é uma subárea da Inteligência Artificial
que desenvolve técnicas computacionais para compreender, interpretar e gerar textos em língua
humana \cite{bird2009}. Engloba um amplo espectro de tarefas, incluindo tokenização,
análise morfossintática, reconhecimento de entidades nomeadas, tradução automática,
sumarização e classificação de textos.

Os avanços dos últimos anos em NLP foram impulsionados por três fatores convergentes:
a disponibilidade de grandes volumes de texto digital, o aumento da capacidade computacional
e o desenvolvimento de arquiteturas de redes neurais profundas capazes de aprender
representações semânticas ricas a partir de dados não rotulados. O paradigma de
pré-treinamento e ajuste fino (\textit{pre-training and fine-tuning}) -- em que um modelo é treinado em
corpus genérico de larga escala e, posteriormente, especializado para uma tarefa específica com
conjunto menor de dados rotulados -- tornou-se dominante na área \cite{devlin2019}.

Para o contexto desta dissertação, as técnicas de NLP mais relevantes são a representação
vetorial de textos (\textit{word embeddings} e \textit{sentence embeddings}), a classificação de textos e,
especialmente, a análise de sentimentos por aspecto, que permite associar polaridades afetivas a
dimensões específicas de um objeto avaliado.

\section{Análise de Sentimentos}
\label{sec:analise-sentimentos}

A análise de sentimentos (\textit{sentiment analysis} ou \textit{opinion mining}) é a tarefa de NLP
que consiste em identificar e extrair automaticamente opiniões, atitudes e emoções expressas em
textos \cite{liu2023b}. Classicamente, opera em três níveis de granularidade: nível de
documento (classificação global do sentimento de um texto), nível de sentença e nível de aspecto
(\textit{aspect-based sentiment analysis} -- ABSA).

A análise por aspecto (\gls{ABSA}) é particularmente relevante para aplicações em políticas
públicas, pois permite identificar \emph{sobre o quê} o sentimento é positivo ou negativo
\cite{wankhade2022}. Por exemplo, em uma avaliação de terminal rodoviário, é possível
discriminar sentimentos relativos à limpeza, à pontualidade dos serviços, à segurança, à
infraestrutura física e ao atendimento dos funcionários, ao invés de atribuir uma polaridade
única ao texto. Essa granularidade é essencial para a geração de diagnósticos acionáveis voltados
a gestores públicos.

As abordagens para análise de sentimentos evoluíram de métodos baseados em léxicos de
polaridade --- como o \gls{VADER} \cite{hutto2014} e o TextBlob \cite{loria2018}, que classificam
textos somando pontuações pré-definidas de palavras --- para métodos baseados em aprendizado
de máquina e, mais recentemente, em modelos neurais profundos. Modelos baseados em
Transformers representam o estado da arte em precisão para análise de sentimento contextual
\cite{wankhade2022}, enquanto \gls{LLM}s generativos emergem como alternativa flexível para
contextos onde dados rotulados são escassos, especialmente para ABSA sem necessidade de
treinamento adicional \cite{agua2025,larosa2025}.

\section{Transformers na NLP}
\label{sec:transformers}

A arquitetura Transformer, introduzida por \citeauthor{vaswani2017}
\citeyear{vaswani2017} com o trabalho ``Attention is All You Need'', representou um marco
fundamental para o Processamento de Linguagem Natural moderno. Diferentemente das
arquiteturas sequenciais como as Redes Neurais Recorrentes (RNNs), os Transformers removeram
a recorrência e fizeram uso intensivo de mecanismos de auto-atenção (\textit{self-attention}),
permitindo o processamento paralelo de \textit{tokens} e uma melhor captura de dependências de
longo alcance no texto. Essa inovação serviu de base para o desenvolvimento de modelos de
linguagem de grande impacto, como \gls{BERT} e \gls{GPT}, que, por sua vez, impulsionaram
avanços significativos em diversas tarefas de NLP, incluindo a análise de sentimentos.

A Revisão Sistemática da Literatura (\gls{RSL}) conduzida como parte deste trabalho
identificou que estudos recentes na área de análise de sentimentos têm se beneficiado dessas
arquiteturas avançadas. Por exemplo, o estudo de \citeauthor{schimanski2024}
\citeyear{schimanski2024} utilizou modelos baseados em Transformers, como RoBERTa e
DistilRoBERTa (versões otimizadas do BERT), para analisar mais de 13,8 milhões de textos
relacionados a comunicações ESG, demonstrando a capacidade dessas técnicas de quantificar
informações com confiabilidade. Concluiu-se que a maioria dos estudos revisados empregou
técnicas avançadas, refletindo a influência dos Transformers no campo de estudo.

\section{Modelo BERT e Suas Aplicações}
\label{sec:bert}

O \gls{BERT} (\textit{Bidirectional Encoder Representations from Transformers}), introduzido
pelo Google em 2018 \cite{devlin2019}, é um dos modelos mais emblemáticos baseados na
arquitetura Transformer. Sua principal característica é a capacidade de considerar o contexto
completo de uma palavra (tanto à esquerda quanto à direita) durante o pré-treinamento,
resultando em representações vetoriais de palavras mais ricas e contextuais. A arquitetura do
BERT é otimizada como um \textit{encoder}, tornando-o particularmente eficaz para tarefas de
classificação de texto, como a análise de sentimentos.

O estudo de \citeauthor{gupta2024} \citeyear{gupta2024} destacou-se ao utilizar o modelo
BERT em conjunto com a técnica YAKE para extração de palavras-chave de relatórios de
sustentabilidade, alcançando 98\% de precisão. Isso confirma a eficácia do BERT na análise de
grandes volumes de dados textuais para obtenção de resultados detalhados.

Modelos BERT, quando submetidos a um processo de \textit{fine-tuning} (ajuste fino) em
\textit{datasets} rotulados para análise de sentimentos, tendem a apresentar alta acurácia (superior a
90\% em \textit{datasets} de referência). Existem também variações e otimizações relevantes: o RoBERTa
apresenta melhor desempenho que o BERT original; o DistilBERT é cerca de 40\% menor e 60\%
mais rápido, mantendo precisão próxima; e o BERTimbau é uma versão pré-treinada para o
português. Embora eficientes, o treinamento e o \textit{fine-tuning} de modelos BERT podem exigir
recursos computacionais significativos e \textit{datasets} específicos rotulados.

\section{Modelos de Linguagem de Grande Porte}
\label{sec:llms}

Modelos de Linguagem de Grande Porte (\gls{LLM}s) genéricos, como as famílias Claude
\cite{anthropic2024}, \gls{GPT} incluindo o GPT-4 \cite{openai2023gpt4}, DeepSeek
\cite{deepseek2024} e Gemini, representam uma evolução significativa no Processamento de
Linguagem Natural, operando de maneira distinta dos modelos como o BERT que são
frequentemente ajustados (\textit{fine-tuned}) para tarefas específicas. Embora um BERT
\textit{fine-tuned} possa ser mais eficaz para uma tarefa de análise de sentimentos para a qual foi
especificamente treinado, os LLMs genéricos oferecem vantagens distintas, especialmente quando
não se dispõe de grandes volumes de dados rotulados ou quando a natureza da tarefa exige maior
flexibilidade e compreensão de nuances.

A utilização de um LLM genérico é particularmente vantajosa em cenários que envolvem
textos complexos com ironia, sarcasmo ou forte contexto cultural (como gírias), pois os LLMs
tendem a performar melhor devido ao seu treinamento em volumes massivos e extremamente
diversificados de dados textuais. Outra situação propícia é a necessidade de aprendizado
\textit{zero-shot} ou \textit{few-shot} \cite{brown2020}: quando não existem dados rotulados para realizar o
\textit{fine-tuning} de um modelo específico, LLMs como Claude ou GPT-4 podem ser empregados
diretamente por meio de \textit{prompts} bem elaborados.

LLMs são igualmente capazes de realizar análise por aspecto (\textit{aspect-based sentiment
analysis}) sem treinamento adicional, identificando sentimentos sobre características específicas
mencionadas no texto. As vantagens estendem-se à flexibilidade analítica: permitem análise
multidimensional por meio de \textit{prompts}, combinando classificação de sentimento com listagem de
razões subjacentes; suportam múltiplas tarefas integradas (análise de sentimento, tradução,
sumarização); e oferecem suporte a contextos longos, com alguns modelos conseguindo analisar
textos com até 200 mil \textit{tokens}, como o Claude 3.5 Sonnet \cite{anthropic2024}.

No entanto, a acurácia dos LLMs pode ser variável e é altamente dependente da qualidade e
clareza do \textit{prompt} fornecido. Existe também risco de ``alucinações'' ou respostas não
padronizadas que não se encaixam em esquemas de classificação rígidos, e o custo computacional
para inferência é geralmente mais alto em comparação com arquiteturas como o BERT. A decisão
entre um BERT \textit{fine-tuned} e um LLM genérico dependerá, portanto, da disponibilidade de dados
rotulados, da complexidade do texto a ser analisado e dos objetivos específicos da análise.
